\chapter{Đánh giá và kiểm thử}

\section{Đối chiếu yêu cầu đề bài}
\begin{longtable}{p{0.38\textwidth}p{0.15\textwidth}p{0.35\textwidth}}
\toprule
\textbf{Yêu cầu} & \textbf{Trạng thái} & \textbf{Bằng chứng} \\
\midrule
Dữ liệu Việt Nam, tối thiểu 7 biến độc lập và 2.000 dòng & \textit{Đạt/Chưa} & \textit{Mục/hình/log} \\
Nguồn đáng tin cậy và xử lý minh bạch & \textit{Đạt/Chưa} & \textit{Mục/hình/log} \\
Biểu đồ phù hợp, rõ ràng, liên kết và tương tác & \textit{Đạt/Chưa} & \textit{Mục/hình/log} \\
Phân tích xu hướng, quan hệ biến và insight chuyên môn & \textit{Đạt/Chưa} & \textit{Mục/hình/log} \\
AI hiển thị code, giải thích và chờ phê duyệt & \textit{Đạt/Chưa} & \textit{Mục/hình/log} \\
Thực thi local và lưu đầy đủ logs & \textit{Đạt/Chưa} & \textit{Mục/hình/log} \\
\bottomrule
\end{longtable}

\section{Kiểm thử chức năng}
Liệt kê ca kiểm thử dashboard, bộ lọc, API dữ liệu, API AI, chỉnh sửa/phê duyệt,
thực thi local, logs và xử lý lỗi.

\section{Kiểm thử an toàn AI}
Trình bày kết quả khi AI sinh code sai, truy vấn ghi dữ liệu, yêu cầu ngoài phạm vi,
code chạy quá lâu hoặc người dùng từ chối phê duyệt.

\section{Đánh giá khả dụng}
Mô tả đối tượng thử nghiệm, nhiệm vụ, phản hồi, vấn đề phát hiện và thay đổi đã thực
hiện. Nếu chưa có nghiên cứu người dùng chính thức, nêu rõ giới hạn.

\section{Hiệu năng và độ ổn định}
Trình bày thời gian tải dashboard, thời gian truy vấn, thời gian phản hồi AI và hành
vi khi API/mạng không khả dụng.
